\cotefile{136}
\scanpage{95}\makeatletter\tagsleft@true\makeatother\mspage{- 169 -}\hfill IX

finies ; donc $X\longmapsto\hat X$ transforme groupes en groupes. Notons que $\hat X$ ne dépend en fait que de la restriction de $X$ à un voisinage ouvert arbitraire $U$ de $S$.

\uline{Lemme} 13.2.1. \uline{Les données sont celles de} 13.1, \uline{et on suppose que les conditions suivantes sont vérifiées} : a) $A=A^{\unsure{o|\circ}}$, b) $R$ \uline{est constant au voisinage de} $s$, \uline{et} c) $R_s$ \uline{engendre l'espace vectoriel complexe} $E_s$. \uline{Soit} $A'$ \uline{un deuxième groupe analytique lisse sur} $S$, \uline{et considérons l'application}
\[
 u\longmapsto\widehat{\unsure{\mathbb{u}|u}}:\operatorname{Hom}_{S\text{-}\mathrm{gr}}(A,A')\longrightarrow\operatorname{Hom}_{\hat S\text{-}\mathrm{gr}}(\widehat{\unsure{\mathbb A|A}},\hat A')\quad,
\]
\uline{où les notations sont celles de} 13.2, \uline{et l'application analogue obtenue en se restreignant au-dessus d'un voisinage ouvert} $U$ \uline{variable de} $s$, \uline{et passant à la limite inductive} :
\[
 u\longmapsto\hat u:\operatorname{Hom}_{(S,s)\text{-}\mathrm{gr}}(A,A')\longrightarrow\operatorname{Hom}_{\hat S\text{-}\mathrm{gr}}(\hat A,\hat A')\quad,
 \tag*{(13.2.1.1)}
\]
(\uline{où} $(S,s)$ \uline{désigne le} germe \uline{d'espace analytique de} $S$ \uline{en} $s$). \uline{Alors l'application précédente est} bijective.

Considérons en effet la suite exacte
\[
 0\longrightarrow R'\longrightarrow E'\longrightarrow {A'}^{o}\longrightarrow\unsure{0|C}
\]
analogue à (13.1.3) définie par $A'$. Compte tenu du fait que les homomorphismes $A\longrightarrow A'$ se factorisent nécessairement par $ {A'}^{o}$, et de même après les changements de base $U\longrightarrow S$ et $S_n\longrightarrow S$, on peut supposer $A'={A'}^{o}$. Compte tenu du dictionnaire 13.1.4, il faut prouver que la donnée d'un
% archive sheet number: 94
% stamp: numbered label 487; 487 also in pencil above the lower strip.
\scanpage{96}
\diagram{\begin{tikzcd}[ampersand replacement=\&,row sep=small,column sep=small]
0\&0\&0\\
0\&0\&\unsure{A|\mathcal A}^{22}\&\text{---}\\
0\&A^{11}\&A^{21}\arrow[u]\&\text{-- --}\\
k\&A^{10}\arrow[u]\&A^{20}\arrow[u]\&\cdots
\end{tikzcd}}

Soit $A^{\unsure{**|\alpha\alpha}}\in\operatorname{Ob}\mathrm{DR}_{\mathrm{sp}}(X,k)$

$T\in\Gamma_X(A^{10})$ définit $\underset{\lambda\unsure{\vdash|\mapsto}\lambda T}{k\longrightarrow A^{1,0}}=\Phi$, \unsure{on a|avec} suite exacte
\diagram{\begin{tikzcd}[ampersand replacement=\&,column sep=small,row sep=small]
0\arrow[r,dash]\&k\arrow[r]\&\Phi\arrow[r]\arrow[d,equal]\&\Psi\arrow[r]\arrow[d]\&0\\
\&\&A^{10}\&A^{20}
\end{tikzcd}}

Comme $\Psi$ est $k$-plat, $\Phi$ \unsure{aussi}, et on sait que
\[
 \mathrm{DR}(\Phi,T)=\underbrace{\Gamma^{\unsure{*|\alpha}}(\Phi)}_{\substack{\text{deg}\\\text{\unsure{compl.|compl}}}}\otimes_k
 \underbrace{\bigwedge^{*}\Psi}_{\text{deg ext}}
\]
muni des \deleted{\illegible{1}}\unsure{structures} \unsure{habituelles} est une $k$-alg. de De Rham \unsure{spéciale} et on a un $k$-hom \unsure{canonique} de telles
\[
 \mathrm{DR}(\Phi,T)\longrightarrow A^{**}
\]
\deleted{donc la cat. de}

Un hom $A^{**}\longrightarrow {A'}^{**}$ définit \unsure{diagramme commutatif}
\diagram{\begin{tikzcd}[ampersand replacement=\&]
\mathrm{DR}(\Phi,T)\arrow[r]\arrow[d]\& A^{**}\arrow[d]\\
\mathrm{DR}(\Phi',T'\unsure{\,|)}\arrow[r]\&{A'}^{**}
\end{tikzcd}}

Cela implique aisément que la cat. \unsure{dérivée} $\mathrm{DR}_{\mathrm{sp}}(X,k)_0$ est \unsure{équivalente canoniquement} à la \deleted{\illegible{1}} catégorie \inserted{\unsure{dérivée|dérivées}} \unsure{déduite|décrite|qui se déduit} des algèbres de la forme $\mathrm{DR}(\Phi,T)$, qui est aussi celle des suites exactes $0\unsure{\longrightarrow|\text{---}} k\longrightarrow\Phi\longrightarrow\Psi\longrightarrow0$ avec $\Psi$ plat sur $k$.
% archive sheet number: 95
\scanpage{97}\mspage{- 156 -}\hfill IX

\[
 d(x_i)=C'_i-C''_i\quad,\tag*{(12.4.2)}
\]
$C'_i$ et $C''_i$ étant les deux composantes irréductibles de $\widetilde C$ qui correspondent aux branches passant par le point double $x_i$. L'opérateur $d$ dépend du choix de l'ordre $(C'_i,C''_i)$ qu'on se donne entre ces deux branches, et l'isomorphisme (12.4.1) dépend évidemment de ce système de choix. Cependant, si pour $i\in I$ on désigne par
\[
 \varphi_{x_i}:M\longrightarrow\mathbb Z\tag*{(12.4.3)}
\]
le composé $M\longrightarrow\mathbb Z^I\xrightarrow{\mathrm{pr}_i}\mathbb Z$, où le premier homomorphisme provient de (12.4.1), on constate aussitôt (en revenant à l'expression (12.3.4)) que $\varphi_{x_i}$ est bien déterminé au signe près ; de façon plus précise, il dépend seulement de l'ordre pris entre les deux branches $C'_i,C''_i$ passant par $i$, et change de signe quand on change cet ordre. Il s'ensuit notamment que la forme billinéaire
\[
 \varphi_{x_i}\otimes\varphi_{x_i}\in\check M\otimes_{\mathbb Z}\check M\simeq\operatorname{Hom}(M\otimes M,\mathbb Z)
 \tag*{(12.4.4)}
\]
est déterminée canoniquement en termes du point singulier $x_i$, sans ambiguité de signe. On peut donc poser
\[
 u=\sum_i\varphi_{x_i}\otimes\varphi_{x_i}:M\otimes_{\mathbb Z}M\longrightarrow\mathbb Z\quad,
 \tag*{(12.4.5)}
\]
et on trouve ainsi un accouplement bien déterminé i.e. indépendant de tout choix de signes. Il est clair qu'il est \uline{défini positif}, car c'est la restriction à $M\hookrightarrow\mathbb Z^I$ de la forme quadratique standart $\sum X_i^2$ sur $\mathbb Z^I$.
% Mirrored arrow below “canoniquement” shows through from the opposite side.
% archive sheet number: 96
% stamp: numbered label 474; 474 also in pencil above the lower strip.
\scanpage{98}
On sait que cette dernière est \unsure{équivalente} à la cat. ponctuelle. Donc $\mathrm{DR}_{\mathrm{sp}}(X,k)_0$ est équivalente à la cat. ponctuelle.

Si maintenant on travaille dans $\mathrm{DR}_{\unsure{\mathrm{spufl}|\mathrm{spu\,fl}}}$ \deleted{(\unsure{resp.|cat.} $\mathrm{DR}_{\unsure{\mathrm{spufl}|\mathrm{spfl}}}$)} i.e. avec des $\unsure{A|\mathcal A}^{\unsure{**|\alpha\alpha}}$ \inserted{u.} flasques, ($A^{ij}$ \inserted{u.} flasque si $0<j$) alors $\Phi$ est \inserted{u.} flasque. \deleted{(\unsure{resp. univ. flasque}).}
Supposons que ceci \unsure{implique} que les $\Gamma_k^{i}\Phi$ $(i>0)$ sont \unsure{univ.} flasques, alors $\mathrm{DR}(\Phi,T)$ est aussi flasque, \deleted{\unsure{et}} et $\unsure{=}$ univ. flasque, et on voit comme dessus que la cat. dérivée est \unsure{équiv.|équivalente} à celle construite sur \deleted{la cat. des} $\mathrm{DR}(\Phi,T)$, avec $\Phi$ univ. flasque et $\Phi/k.T$ plat, i.e. avec les \unsure{suite|suites} exactes $0\unsure{\longrightarrow|\text{---}} k\longrightarrow\Phi\longrightarrow\Psi\longrightarrow0$, avec $\Phi$ univ. flasque et $\Psi$ $k$-plat ; or cette cat. dérivée est encore la catégorie ponctuelle. On trouve donc

\uline{Prop} La cat. $\mathrm{DR}_{\mathrm{sp}}(X,k)_0$ est équivalente à la cat. ponctuelle. Il en est de $=$ de la cat. $\mathrm{DR}_{\mathrm{sp\,ufl}}(X,k)_0$, si on suppose que pour \unsure{tt} \inserted{$k$-Module} $\Phi$ plat et univ. flasque, les $\Gamma^i\Phi$ $(i>0)$ sont aussi univ. flasques.
% archive sheet number: 97
\scanpage{99}\mspage{- 167 -}\hfill IX

$A$ \uline{lisse} sur $S$, et on désigne par
\[
 \underline E=\underline{\mathrm{Lie}}(A/S)\tag*{(13.1.1)}
\]
le module localement libre sur $S$, image inverse par la section nulle du faisceau tangent relatif de $A$ sur $S$. Nous admettrons ici la définition de l'homomorphisme exponentiel
\[
 \exp:E\longrightarrow A\quad,\tag*{(13.1.2)}
\]
où $E$ est le fibré vectoriel analytique sur $S$ dont le Module des sections holomorphes est $\underline E$. Cet homomorphisme se définit, au voisinage d'un point $s\in S$, d'abord au voisinage de la section nulle de $E$ par les méthodes bien connues en théorie des groupes formels en car. nulle, puis se prolonge de façon évidente par multiplicativité à $E$ tout entier. C'est un homomorphisme de groupes analytiques relatifs sur $S$, qui est \uline{étale} et a comme image le sous-espace ouvert $A^{\unsure{o|\circ}}$ de $A$ réunion des composantes neutres des fibres de $A$ sur $S$. Si $R$ est le sous-espace analytique noyau de $\exp$, on a donc une suite exacte de groupes analytiques relatifs :
\[
 0\longrightarrow R\longrightarrow E\longrightarrow A^o\longrightarrow0\quad,\tag*{(13.1.3)}
\]
où le morphisme $E\longrightarrow A^o$ est surjectif et étale, donc $R$ est un groupe analytique relatif \uline{étale} sur $S$. C'est un sous-espace analytique de $E$, qui est fermé dans $E$ si et seulement si $A^o$ est séparé sur $S$, i.e. la section nulle de $A^o$ est fermée dans $A^o$.
% archive sheet number: 98
% stamp: numbered label 485; 485 also in pencil above the lower strip.
\scanpage{100}
Cette dernière condition est satisfaite pour le topos \inserted{$(\mathrm{Ss})=$} $\widehat\Delta$, car pour celui-ci
% slanted
\margin{faux tel quel, à arranger!}
$\underbrace{\text{flasque}+\underset{\downarrow}{\text{plat}}}_{\Downarrow}$ \inserted{\uline{\unsure{non|univ.}, projectif}} $\Longrightarrow$ homotope à $0$, donc univ. flasque et les $\Gamma^i(\Phi)$ univ. flasques. $(i>0)$.

Soit $P$ le topos ponctuel, et
\[
 i_n:P\longrightarrow\unsure{\widehat\Delta|\widehat{\mathbb\Delta}}(\mathrm{Ss})
\]
le morphisme de topos associé au foncteur fibre
\[
 u_n^*(K)=K_n=K(\Delta_n).
\]
Soit $k$ un anneau (\unsure{constant|commutatif}).

\uline{Lemme} $\forall$ $k$-module $M$, on a un isom. can.
\[
 i_{0*}(M)\simeq\Phi_*\otimes_{\mathbb Z}M\simeq\Phi_{*k}\otimes_k M
\]
où $\Phi_*\in\operatorname{Ob}(\widehat\Delta_{\mathrm{ab}}\simeq\underline{\mathrm{Hom}}(\Delta^o,\mathrm{Ab}))$ est défini par $\Phi_{\unsure{*|k}}(I)=\mathbb Z^I$, \unsure{d'où|donc} $\Phi_{*k}(I)\simeq k^I$.

En effet, on a par construction de l'image directe
\[
\begin{aligned}
i_{0*}(M)(I)&\simeq M(i_0^*(I))\simeq M^I\\[-2pt]
\rotatebox{90}{$=$}\quad&\qquad\rotatebox{90}{$=$}\\[-2pt]
\operatorname{Hom}(I,i_{0*}(M))&\qquad\operatorname{Hom}(\underbrace{i_0^*(I)}_{\operatorname{Hom}(\Delta_0,I)\simeq I},M).
\end{aligned}
\]
OK.

\deleted{Soit \unsure{$k$.}}

\uline{Corollaire} Pour tout $k_*$-Module $\deleted{\unsure{C|E}}A_*$, on a
\[
 \operatorname{Hom}_{k_*}(\deleted{\unsure{C}}A_*,\Phi_*\otimes_{\unsure{k|\mathbb Z}}M)\simeq\operatorname{Hom}_k(\deleted{\unsure{C}}A_0,M)
\]
en particulier on a un hom can. (pour $M=C_0$)
\[
 A_*\longrightarrow\Phi_*\otimes_{\mathbb Z}\deleted{\unsure{M}}A_0\simeq\Phi_{*k}\otimes_k A_0\quad(=i_{0*}i_0^*A_*)
\]
\deleted{\unsure{Nous|Hom}} Appliquons ceci au cas d'une immersion
\[
 k_*\overset{u}{\hookrightarrow}\deleted{\unsure{C}}A_*=\deleted{\unsure{C}}A_*^{1,0}
\]
provenant d'une $\unsure{A|\mathcal A}_*^{\unsure{**|\alpha\alpha}}\in\operatorname{Ob}\mathrm{DR}_{\mathrm{sp}}(\mathrm{Ss},k_*)$
% archive sheet number: 99
\scanpage{101}\mspage{- 171 -}\hfill IX

engendre $R$ au voisinage de $s$, d'où il résulte que $F(R)\subset R'$ au voisinage de $s$, ce qui achève la démonstration.

\uline{Corollaire} 13.2.2. \uline{Si} $A$ \uline{satisfait les conditions} a), b), c) \uline{de} 13.2.2, \uline{il est déterminé à isomorphisme unique près au voisinage de} $s$ \uline{par la connaissance du} $\hat S$-\uline{groupe} $\hat A$.

13.3. \uline{L'extension de \unsure{Raynaud} analytique}. Nous sommes maintenant en mesure de paraphraser, en Géométrie Analytique, la construction de l'extension de Raynaud donnée dans 7 . Supposons donné $A$ comme dans 13.1, avec $A=A_{\unsure{o|0}}$, et considérons $R_o=R_s$, qui est un $\mathbb Z$-module libre de type fini. Quitte à restreindre $S$ à un voisinage ouvert assez petit de $s$, interprétant $R_o$ comme le groupe des germes de sections de $R$ au voisinage de $s$, on trouve un homomorphisme
% Short unattached horizontal stroke above the space in « trouve un ».
\[
 R_{oS}\longrightarrow R\subset E \tag*{(13.3.1)}
\]
où $R_{oS}$ désigne le groupe analytique constant sur $S$ de valeur $R_o$. Choisissant le voisinage de $s$ assez petit pour que toute section de $\underline O_S$ dont le germe en $s$ est nul soit nul, on voit que l'homomorphisme précédent est injectif. Son image est donc un sous-groupe analytique constant de $R$, qui est évidemment le plus grand sous-groupe constant de $R$, et qui s'appelle la partie \uline{fixe} de $R$, et est notée $R^f$ :
% Short unattached horizontal stroke below the space in « un sous-groupe ».
\[
 R_{oS}\simeq R^f\lhook\joinrel\relbar R\,,\ \text{induit un isomorphisme } R^f_s\simeq R_s\,.
 \tag*{(13.3.2)}
\]
On pourra alors considérer le groupe quotient
\[
 A^{\handwritten{\natural}}=E/R^f\tag*{(13.3.3)}
\]
% archive sheet number: 100
% stamp: numbered label 489; faint pencil numerals above the lower strip.
\scanpage{102}
\deleted{Considérons l'hom induit en degré $0$}
\[
 \deleted{k_0=k\longrightarrow A_0}
\]
On trouve alors un hom \inserted{donnant lieu à un diag. \unsure{comm|comm.}}
\diagram{\begin{tikzcd}[ampersand replacement=\&]
\deleted{\unsure{C|E}}A_*\arrow[r]\&\Phi_*\otimes_{\mathbb Z}\deleted{\unsure{C}}A_0\\
k_*\arrow[u,hook]\arrow[r,hook]\&\Phi_{*k}(=\Phi_*\otimes_{\mathbb Z}\unsure{k_0|k})\arrow[u,hook,swap,"{\Phi_*\otimes_{\mathbb Z}u_0}"]
\end{tikzcd}}

Ainsi on trouve un diagramme \unsure{canonique} d'hom dans $\mathrm{DR}_{\mathrm{sp}}(\mathrm{Ss},k_*)$
\diagram{\begin{tikzcd}[ampersand replacement=\&,column sep=small]
\mathrm{DR}(A_*^{10},T)\arrow[r]\arrow[d]\&A_*^{\unsure{**|\alpha\alpha}}\\
\mathrm{DR}(A_0^{10}\otimes_{\mathbb Z}\Phi_*,T)\&\mathrm{DR}(\Phi_{*k},\overset{\nearrow}{T})\arrow[l]
\end{tikzcd}}
reliant une \inserted{$k$-}algèbre de DR. \unsure{spéciale}, en 3 crans, à l'algèbre de DR \unsure{standart|standard} $\mathrm{DR}(\Phi_{*k},T)$ (qui semble la plus économique de \unsure{toute|toutes}). Si $A_*^{**}$ est \inserted{(\unsure{univ})} flasque, alors il en est de $=$ de $\mathrm{DR}(A_*^{10},T)$, et le diagramme ci-dessus se trouve dans $\mathrm{DR}_{\unsure{\mathrm{spUfl}|\mathrm{spufl}|\mathrm{spfl}}}$.

On va introduire maintenant les catégories $\mathrm{DR}_{\mathrm{sp}}'(X,k)$, $\mathrm{DR}_{\mathrm{sp\,ufl}}'(X,k)$ ayant les $=\mathrm{s}$ objets que les catégories sans $'$, mais ayant plus de morphismes : les hom. de $k$-algèbres \deleted{\unsure{catég.}}\inserted{différentielles} \deleted{à \unsure{valeurs}} à p.d., respectant $T$
% archive sheet number: 101
\scanpage{103}\mspage{- 166 -}\hfill IX

qui n'est autre que 12.5 dans le cas de $\underline X^\lambda$ sur $\underline S^\lambda$. Or $\underline S$ est un trait de \uline{caractéristique nulle}. Cela achève donc la réduction annoncée.

\uline{Remarques} 12.9. On notera que cette réduction nous réduit en même temps au cas où $X_{\unsure{o|0}}$ n'a qu'un seul point singulier, qui est celui qui était considéré classiquement dans la formule de Picard-Lefschetz. Notons d'ailleurs que dans ce cas, $X_o$ a une ou deux composantes irréductibles, ces cas correspondants respectivement à $\operatorname{rang} M=1$ et $\operatorname{rang} M=0$ (ce dernier cas étant celui où la jacobienne de la fibre générique a bonne réduction).

12.10. La démonstration donnée dans 12.8 nous fournit un résultat plus général que 12.5, valable sans supposer que $X$ soit régulier, en introduisant les entiers $d_\lambda\geq1$ de (12.8.12) : on trouve alors que les accouplements de monodormie se déduisent tous de l'accouplement
\[
 u=\sum_{\lambda\in L}d_\lambda\varphi_{x_\lambda}\otimes\varphi_{x_\lambda}:M\unsure{\otimes|\hat\otimes}M\longrightarrow\mathbb Z\quad.
 \tag*{(12.10.1)}
\]

13. \uline{Liens avec la théorie transcendante} : \uline{cas analytique complexe}.

Pour les notions générales relatives aux espaces analytiques, nous renvoyons à [2].

13.1. \uline{Rappels sur les groupes analytiques sur} $S$. Soit $S$ un espace analytique, $A$ un espace analytique au-dessus de $S$, muni d'une structure de groupe commutatif au-dessus de $S$ (i.e. $A$ est un groupe commutatif de la catégorie des espaces analytiques au-dessus de $S$). On suppose
% archive sheet number: 102
% stamp: numbered label 484; 484 also in pencil above the lower strip.
\scanpage{104}
\[
 u:\underset{\rotatebox{90}{$=$}}{\widehat{\unsure{A|\mathcal A}}^{\unsure{**|\alpha\alpha}}}\mathrel{\unsure{\relbar|\longrightarrow}}
 \underset{\rotatebox{90}{$=$}}{{\widehat A'}^{**}}
\]
\[
 \prod A^{ij}\qquad\prod {A'}^{ij}
\]
hom d'Anneaux \unsure{diff|de}, appliquant $A^{**+}$ dans ${A'}^{**+}$ et compatible aux p.d ; et tels que $u(T)=T$ (NB a priori
\[
 u:\underset{H^{*}(A^{**})}{\underset{\rotatebox{90}{$\simeq$}}{k\{T\}}}
 \longrightarrow
 \underset{H^*({A'}^{**})}{\underset{\rotatebox{90}{$\simeq$}}{k\{T\}}},
\]
hom compatible aux p.d.). \deleted{À vrai dire, il serait plus intéressant de \unsure{prendre}}

On désigne par
\[
 \mathrm{DR}_{\mathrm{sp}}'(X,k)_0,\quad \mathrm{DR}_{\mathrm{sp\,ufl}}'(X,k)_0
\]
les cat. dérivées associées (i.e. les groupoïdes engendrés). On a donc un \unsure{diagramme|\deleted{diagramme}|diagrammes} comm. de foncteurs \unsure{can}
\diagram{\begin{tikzcd}[ampersand replacement=\&,column sep=large,row sep=huge,cells={font=\scriptsize}]
\mathrm{DR}_{\mathrm{sp}}(X,k)
 \arrow[r,dash,"{\text{surj. sur objet}}","{\unsure{\cdot|\rightarrow}}"{pos=.98,inner sep=0pt}]
 \arrow[d,"{\substack{\deleted{\unsure{\text{iso.}}}\ \text{surjectif}\\\text{sur obj}}}"]
 \&\mathrm{DR}_{\mathrm{sp}}(X,k)_0\overset{\sim}{\simeq}\text{ponctuelle}\arrow[d]\\
\mathrm{DR}_{\mathrm{sp}}'(X,k)\arrow[r,swap,"{\substack{\text{surj. sur}\\\text{objet}}}"]
 \&\mathrm{DR}_{\mathrm{sp}}'(X,k)_0\\
\mathrm{DR}_{\mathrm{sp\,ufl}}(X,k)\arrow[uu,bend left=12,"{\deleted{\unsure{\blacktriangledown}}}"{pos=.07,inner sep=0pt}]
 \arrow[d,"{\substack{\deleted{\unsure{\text{iso}}}\ \text{surjectif}\\\text{sur obj}}}"]
 \arrow[r,"{\substack{\text{surj.}\\\text{sur obj}}}"]
 \&\mathrm{DR}_{\mathrm{sp\,ufl}}(X,k)_0\overset{\sim}{\simeq}\text{ponct.}
 \arrow[uu,bend left=12]\arrow[d,"{\substack{\text{(si ufl.}\\\text{stable par}\\\text{les }\Gamma^i,\ i>0)}}"]\\
\mathrm{DR}_{\mathrm{sp\,ufl}}'(X,k)\arrow[uu,bend left=13,"{\deleted{\unsure{\blacktriangledown}}}"{pos=.07,inner sep=0pt}]
 \arrow[r,"{\substack{\text{surj.}\\\text{sur}}}","{\text{obj}}"']
 \&\mathrm{DR}_{\mathrm{sp\,ufl}}'(X,k)_0\unsure{\simeq|=}\arrow[uu,bend left=13]
\end{tikzcd}}

qui montre que les deux catégories \inserted{\unsure{dér.}} $\mathrm{DR}_{\mathrm{sp}}'(X,k)_0$ et $\mathrm{DR}_{\mathrm{sp\,ufl}}'(X,k)_0$ sont des groupoïdes \inserted{\unsure{trs|ens|, en}} \unsure{pointés|ponctués|fait triviaux} (munis d'une famille d'objets avec syst. transitif d'iso. entre). Donc ils sont can
% archive sheet number: 103
\scanpage{105}\mspage{- 170 -}\hfill IX

germe d'homomorphisme de fibrés vectoriels $E\longrightarrow E'$, appliquant $R$ dans $R'$, revient à la donnée d'un système cohérent d'homomorphismes $E_n\longrightarrow E'_n$ appliquant $R_n$ dans $R'_n$. Or si $\underline O_s$ désigne l'anneau local de $S$ en $s$, la donnée d'un système cohérent d'homomorphismes $E_n\longrightarrow E'_n$ revient à la donnée d'un homomorphisme de $\widehat{\underline O}_s$-modules
\[
 \hat F:\widehat{\underline E}_s\longrightarrow\widehat{\underline E}'_s
\]
(où $\underline E_s$ et $E'_s$ sont les $\underline O_s$-modules libres fibres en $s$ des $\underline O_s$-Modules $\underline E,\underline E'$ correspondant à $E,E'$), tandis qu'un germe d'homomorphisme de fibrés vectoriels $E\longrightarrow E'$ revient à la donnée d'un homomorphisme de $\underline O_s$-modules
\[
 F:\underline E_s\longrightarrow\underline E'_s\quad.
\]
Cela montre déjà que l'application (13.2.1.1) est injective, et pour la surjectivité, il suffit de prouver que si on a un homomorphisme $\hat F$ comme dessus, tel que pour tout $n$ l'homomorphisme induit $F_n:E_n\longrightarrow E'_n$ applique $R_n$ dans $R'_n$, alors $\hat F$ provient d'un $F$ i.e. applique $\underline E_s$ dans $\underline E'_s$, et le germe correspondant d'homomorphismes $E\longrightarrow E'$ applique $R$ dans $R'$. Or, comme par hypothèse c), $R_s=R_{\unsure{o|0}}$ engendre $E_s$, il s'ensuit par Nakayama que le groupe $R_o$ (identifié au groupe des sections de $\hat R$ sur $\hat S$, ou,au choix,au groupe des germes de sections de $R$ sur $S$) engendre $\hat E_s$ ; donc, pour prouver que $\hat F(\underline E_s)\subset\underline E'_s$,il suffit de prouver que $\hat F(R_o)\subset\underline E'_s$. Or on a $\hat F(R_o)\subset R'_o$, où $R'_o$ est lui-même interprété comme le groupe des germes de sections de $R'\subset E'$ au voisinage de $s$, d'où l'inclusion annoncée, et par suite l'existence de $F$. De plus on voit que $F(R_o)\subset R'_o$, où $R_o$ et $R'_o$ sont considérés comme des groupes de germes de sections de $E$ resp. $E'$. En vertu de b), $R_o$

\diagram{\begin{tikzcd}[ampersand replacement=\&,row sep=small]
\handwritten{\mathrm{DR}(\unsure{A|\mathcal A}_{\unsure{*|\,}}^{1,0},T)}\arrow[d]\arrow[r]\&\handwritten{A^{\unsure{**|\alpha\alpha}}}\\
\handwritten{\mathrm{DR}(\deleted{\unsure{A}}\Phi_*\otimes_{\mathbb Z}A_0,T)}\\
{}\arrow[u]
\end{tikzcd}}
% archive sheet number: 104
% stamp: numbered label 488; 488 also in pencil above the lower strip.
\scanpage{106}
\unsure{équivalents|équivalent} \unsure{à des|à deux|au} groupoïdes \unsure{connexes|connexe} définis par deux groupes, que nous nous proposons de déterminer [dans le cas $X=(\mathrm{Ss})$.]

\diagram{\begin{tikzcd}[ampersand replacement=\&,column sep=small]
\&\widehat{\unsure{A|\mathcal A}}_*^{\unsure{**|\alpha\alpha}}\arrow[r,"u"]\&{\widehat A'}_*^{**}\\
\&\widehat{\mathrm{DR}}(\underset{(A_*)}{\underset{\rotatebox{90}{$=$}}{A_*^{10}}},T)
 \arrow[u]\arrow[ur,dashed,"{\unsure{v|\nu|5|\exists}}"]\\
\widehat{\Gamma^*(A_*)}\mathbin{\widehat\otimes}_k\underset{A_*/k_*}{\underset{\rotatebox{90}{$=$}}{\widehat{\bigwedge^*(B_*)}}}\arrow[ur,equal]\&\&{}
\end{tikzcd}}
\[
 \begin{aligned}
 A_*&\xrightarrow{v}{A'}^{**+}\\
 T&\unsure{\vdash|\longmapsto} T
 \end{aligned}
\]
\[
 v(p,q): A_*\longrightarrow {A'}^{pq}\qquad(p,q\geq0,\ p+q>0)
\]
\[
 \begin{aligned}
 v(1,0)(T)&=T\\
 v(p,q)(T)&=0\quad\text{si }(p,q)\ne(1,0)\\
 \bar v(p,q)&:\deleted{\unsure{\Phi}}B\longrightarrow {A'}^{pq}\quad(\text{si }\unsure{(p+q)|(p,q)}\ne(1,0))
 \end{aligned}
\]
\diagram{\begin{tikzcd}[ampersand replacement=\&,row sep=small,column sep=small]
k\arrow[r,phantom,"{\lhook\joinrel\relbar}" description]\arrow[d,equal]\&A\arrow[d,"{v(1,0)}"]\\
k\arrow[r,phantom,"{\lhook\joinrel\relbar}" description]\&{A'}^{1,0}
\end{tikzcd}}
\deleted{$k\lhook\joinrel\relbar A$}
\diagram{\begin{tikzcd}[ampersand replacement=\&,row sep=small,column sep=small]
\&\&\Phi_*\otimes A\\
k\arrow[r,phantom,"{\lhook\joinrel\relbar}" description]\arrow[drr,dash]\&A\arrow[ur,dash]\arrow[r]\arrow[dr,dash]\&B\arrow[d]\\
\&\& {A'}^{pq}=
\end{tikzcd}}
% archive sheet number: 105
\scanpage{107}\mspage{- 168 -}\hfill IX

13.1.4. Bien entendu, la suite exacte (13.1.3) est \uline{fonctorielle en} $A$, et commute à tout changement de base $S'\longrightarrow S$. De cette façon on trouve une équivalence de catégories, compatible aux changements de base, entre la catégorie des groupes analytiques relatifs sur $S$ lisses et à fibres connexes, et la catégorie des couples $(E,R)$ d'un fibré vectoriel localement libre $E$ sur $S$ et d'un sous-groupe analytique $R$ de $E$ étale sur $S$.

13.1.5. Lorsque $A$ est un groupe ``\uline{abéloide}'' sur $S$, par quoi nous entendons qu'il est lisse, propre et séparé sur $S$ et à fibres connexes, de sorte qu'on a en particulier $A=A^{\unsure{o|\circ}}$, alors $R$ est un groupe localement constant sur $S$, dont le rang en chaque point $s$ (en tant que $\mathbb Z$-module) est égal à $2\operatorname{rang}_s(E)$, et réciproquement. De cette façon on trouve l'interprétation bien connue d'un espace abéloide sur $S$ en termes d'un ``\uline{système local}'' $R$ sur $S$ et d'une structure de fibré vectoriel analytique \uline{complexe} sur le fibré analytique réel $R\otimes_{\mathbb Z}\underline{\mathbb R}$ défini par $R$.

13.2. \uline{Complétion formelle le long d'une fibre}. Nous supposons donné désormais un point $s\in S$, et allons travailler localement au voisinage de $s$. Pour tout entier $n\geq0$, désignons par $S_n$ le $n$.ième voisinage infinitésimal de $s$ dans $S$, et pour tout espace analytique $X$ sur $S$, soit de même $X_n=X\times_S S_n$ ; pour $n$ variable, le système des $X_n$ forme alors un système inductif d'espaces analytiques sur $S$, qu'on dénotera par $\hat X$ et appelera le \uline{complété formel de} $X$ \uline{le long de} $X_s$. Il dépend évidemment fonctoriellement de $X$, et sa formation commute aux limites projectives
% The leftward arrow above “catégorie des couples” shows through from the other side.
% archive sheet number: 106
% stamp: numbered label 486; 486 also in pencil above the lower strip.
\scanpage{108}
Détermination des hom. \deleted{\unsure{dans}} \deleted{$\mathrm{DR}_{\mathrm{sp}}$} de $\unsure{\widehat{\mathrm{DR}}_{\mathrm{st},k}^{**}|\widehat{\mathrm{DR}}_{*k}^{**}|\mathrm{DR}_{*k}^{**}}=
\widehat{\unsure{\Gamma_k^*(\Phi_{*k})\otimes_k\bigwedge_k^*\Psi_{*k}|\Gamma_k(\mathfrak G_{*k})\mathbin{\widehat\otimes}\widehat\Lambda_k(\mathfrak G_{*k}/k)}}$ dans une $k$-Algèbre de de Rham \unsure{spéciale}, \unsure{pour les} seules structures de $k$-Alg \unsure{augmentées}, à p.d., \unsure{munies d'une différentielle}, et à élém\textsuperscript{t} marqué $T$.

C'est \unsure{défini} par $\unsure{\varphi}:\Phi_{*k}\longrightarrow\widehat{\unsure{A|\mathcal A}}_*^{\unsure{**+|\alpha\alpha+}}$, \unsure{$k$-lin}, sous \unsure{conditions} \inserted{\unsure{suiv.}}, que son image \deleted{\unsure{est}} soit \unsure{commutative}, \unsure{de $=$ $d$ de celle-ci}, et que $d$ soit \unsure{nulle} sur \illegible{1} et \unsure{qu'elle} \illegible{2} \unsure{commutent}, et enfin que $\varphi(T)=T$.

Or $\Phi_{\unsure{k|*}}$ est \unsure{comme engendré} \unsure{par les} \illegible{2} du \unsure{degré $1$} : \unsure{la donnée de} $\unsure{A_0^{**+}|\widehat A^{**+}}$ \illegible{2} \unsure{et de} $\unsure{T}$, et $\varphi_1:\underset{\unsure{k\delta^0+k\delta^1}}{\underset{\rotatebox{90}{$=$}}{k\times k}}\longrightarrow\widehat A_1^{**+}$, \unsure{appliquant} \unsure{l'élément} \unsure{sur} $T$ donc il suffit de connaître l'image
$\unsure{\zeta_0|\xi_0|\gamma_0}=\varphi_1(\unsure{\delta^0})$, avec
\[
 \zeta_0\in\widehat A_1^{**+}
\]
et compatible avec les deux \unsure{augmentations} $\alpha^*$ et $\beta^*:\widehat A_1^{**+}\longrightarrow\widehat A_0^{**+}$
\[
 \boxed{\alpha^*(\zeta_0)=1,\quad\beta^*(\unsure{\zeta_0|\xi_0})=\unsure{1|0|T}}\qquad
 (\alpha,\beta:\Delta_0\longrightarrow\Delta_1)
\]
les conditions de \unsure{commutativité} \deleted{\unsure{comm.}}\inserted{\unsure{mentionnées}} \unsure{plus haut} s'explicitent alors \unsure{par}
\[
 \boxed{\unsure{\zeta_0\,d\zeta_0=d\zeta_0\,\zeta_0|\xi_0\,d\xi_0=d\xi_0\cdot\xi_0}}\qquad \boxed{\unsure{d\zeta_0\cdot d\zeta_0=0}}
\]
et si $q_0,q_1$ sont les deux hom \unsure{surj.} $\Delta_2\longrightarrow\Delta_1$ ($q_0(0)=0$, $q_0(1)=q_0(2)=1$, $q_1(0)=q_1(1)=0$, $q_1(\unsure{2})=1$), \unsure{par|puis}
\[
 \boxed{\begin{aligned}
 q_0^*(\zeta_0)q_1^*(\zeta_0)&=\unsure{q_1^*(\zeta_0)q_0^*(\zeta_0)|q_1^*(\xi_0)q_0^*(\xi_0)},\qquad
 \unsure{q_0^*(\zeta_0)\cdot\deleted{\unsure{d}}q_1^*(d\zeta_0)=\deleted{\unsure{d}}q_1^*(d\zeta_0)\cdot q_0^*\zeta_0}\\
 &\hspace{1em}\unsure{q_1^*(\zeta_0)\cdot\deleted{\unsure{d}}q_0^*(d\zeta_0)=\deleted{\unsure{d}}q_0^*(d\zeta_0)\cdot q_1^*\zeta_0}
 \end{aligned}}
\]
\[
 \boxed{\unsure{q_0^*(d\zeta_0)q_1^*(d\zeta_0)=-q_1^*(d\zeta_0)q_0^*(d\zeta_0)|q_0^*(d\xi_0)q_1^*(d\xi_0)=-q_1^*(d\xi_0)q_0^*(d\xi_0)}}
\]
% archive sheet number: 107
\scanpage{109}
\mspage{N27}
Thus the subalgebra $\Lambda_{1,n}$ of the model precisely determines the nilpotent quotient $\pi_1/C_n$ tensored by $Q$ or $R$.

It is perhaps important to note that the minimal model $\Lambda$ of the de Rham algebra is constructed in a purely algebriac computational construction as indicated in the text above. The structure of the model then determines the form of the lower central series of $\pi_1$ (over $Q$ and $R$) and so on for\unsure{:} the further rational or real homotopy theory.

The subalgebra $\Lambda_1$ of $\Lambda$ mapping to the de Rham algebra is constructed \unsure{'}so we have an isomorphism of $H^1$ and an injection on $H^2$. Then 2 dim'l generators are added in a possibly infinite sequence of steps. The first step creates an isomorphism on $H^2$ and the later steps produce an injection on $H^3$. Then 3 dim'l generators are added in a possibly infinite sequence of steps, 4 dim'l, and so on to build the complete model.

The structure of model over $Q$ for the rational de Rham algebra of a space corresponds (via the Hirsch calculation described in Appendix H) to a tower of spaces and maps. Each map in the tower is a principal fibration where the fibre has one non-zero homotopy group which is a rational vector space.

This tower of spaces is the rational homotopy type of the space which is defined in general for all spaces regardless of the fundamental group.

We have described above the relationship between the $\Lambda_1$ of the model (and therefore the tower) and the lower central\unsure{:} series of the fundamental group of the space. The significance of the higher dimensional structure of the model and the tower for the general non-simply connected spaces\unsure{\textperiodcentered} is somewhat complicated and not yet completely understood except in a philosophical way.
% archive sheet number: 108
\scanpage{110}
% Blank ruled/perforated verso; the corner number is archival.
% archive sheet number: 109
